\documentclass[11pt]{article}
\usepackage[a4paper,margin=2.2cm]{geometry}
\usepackage{xcolor}
\usepackage{fontspec}
\usepackage{xeCJK}
\usepackage{enumitem}
\setmainfont{Times New Roman}
\setCJKmainfont{SimSun}
\setlist[itemize]{leftmargin=1.4em,itemsep=0.35em,topsep=0.35em}
\setlength{\parindent}{0pt}
\setlength{\parskip}{0.55em}
\pagestyle{plain}
\begin{document}
{\color[HTML]{1F5FD2}\LARGE\bfseries 感觉统合训练活动手册\par}
{\color[HTML]{667085}\small 星轨原创教育资源：用于家庭与学校共同制定支持计划。\par}
\vspace{0.8em}
{\color[HTML]{111827}\large\bfseries 安全原则\par}
\begin{itemize}
\item 感觉统合活动的目的在于支持学习和日常参与，不应为了追求强刺激而增加儿童负担。
\item 实施前要确认场地安全、儿童意愿、医疗限制和成人看护，出现不适时立即停止。
\end{itemize}
{\color[HTML]{111827}\large\bfseries 本体觉活动\par}
\begin{itemize}
\item 推墙、动物走、搬软垫、拉弹力带和运送较重的课堂材料，都可以帮助儿童感受身体位置和力量。
\item 每轮活动以三到五分钟为宜，观察儿童是否疲劳、烦躁或过度兴奋，再决定是否继续。
\end{itemize}
{\color[HTML]{111827}\large\bfseries 前庭与触觉活动\par}
\begin{itemize}
\item 平衡路径、轻柔摇摆、跳跃游戏、触觉箱、黏土和布料配对应循序渐进，避免突然提高难度。
\item 对触觉敏感的儿童可先用工具接触材料，再过渡到短时间手部接触。
\end{itemize}
\vfill
{\color[HTML]{667085}\footnotesize 本材料为应用内原创教育内容，不能替代医疗、法律或正式评估建议。\par}
\end{document}
